% !TEX program = xelatex
% ============================================================
%  Arabic Scientific Poster Template (Landscape) — Nibras Center
%  Compile with: xelatex main.tex
%  Uses article class + tikz for a single-page A0 landscape poster.
% ============================================================
%  To change poster size, modify the geometry package:
%    A0 landscape: a0paper,landscape  (1189x841mm)
%    A1 landscape: a1paper,landscape  (841x594mm)
% ============================================================

\documentclass[11pt]{article}

% ---- Poster geometry (A0 landscape) ----
\usepackage[a0paper,landscape,margin=2cm,top=2cm,bottom=2cm]{geometry}

% ---- Core packages ----
\usepackage{amsmath}
\usepackage{amssymb}
\usepackage{tikz}
\usepackage{enumitem}
\usepackage{xcolor}
\usepackage{fontspec}

% ---- RTL / Arabic support ----
\usepackage{bidi}
\usepackage[no-math]{fontspec}
\usepackage[quiet]{polyglossia}

% ---- Fonts (explicit sizes via \fontsize, no global Scale) ----
\setmainfont[Scale=1]{NotoKufiArabic-Regular.ttf}
\newfontfamily\arabicfont[Script=Arabic,Scale=1]{NotoKufiArabic-Regular.ttf}
\newfontfamily\englishfont[Scale=1]{TeX Gyre Termes}
\setmonofont{NotoKufiArabic-Regular.ttf}

% ---- Languages ----
\setdefaultlanguage[calendar=gregorian,hijricorrection=1]{arabic}
\setotherlanguage[variant=british]{english}

% ---- Hyperref (load last) ----
\usepackage[hidelinks]{hyperref}

% ---- Colors (customize to your institution) ----
\definecolor{primaryColor}{HTML}{1A3C6B}
\definecolor{secondaryColor}{HTML}{2E6DA4}
\definecolor{lightBg}{HTML}{F0F4F8}
\definecolor{accentColor}{HTML}{D4A017}

% ---- Custom block command (colored box for poster sections) ----
\newcommand{\posterblock}[2]{%
  \begin{tikzpicture}
    \node[fill=lightBg, draw=primaryColor, line width=2pt, rounded corners=5pt,
          inner sep=1cm, text width=\dimexpr\linewidth-2.5cm\relax,
          align=right] {%
      \begin{minipage}{\dimexpr\linewidth-2.5cm\relax}
        {\fontsize{20}{24}\selectfont\bfseries\color{primaryColor} #1}\\[0.25cm]
        {\fontsize{13}{16}\selectfont #2}
      \end{minipage}
    };
  \end{tikzpicture}
  \vspace{0.6cm}
}

% ---- List spacing for poster ----
\setlist[itemize]{topsep=0.4em, itemsep=0.2em, parsep=0pt, leftmargin=1.5cm}
\setlist[enumerate]{topsep=0.4em, itemsep=0.2em, parsep=0pt, leftmargin=1.5cm}

\begin{document}
\pagestyle{empty}

% ============================================================
% TITLE BANNER (full width, landscape)
% ============================================================
\begin{tikzpicture}
  \fill[primaryColor] (0,0) rectangle (\textwidth, -5cm);
  % Title — centered, white, large
  \node[white, align=center, text width=0.85\textwidth] at (\textwidth/2, -1.7cm) {
    {\fontsize{42}{52}\selectfont\bfseries عنوان البحث العلمي}
  };
  % Authors / affiliation — centered, white, medium
  \node[white, align=center, text width=0.85\textwidth] at (\textwidth/2, -3.6cm) {
    {\fontsize{22}{28}\selectfont د. اسم الباحث \quad|\quad القسم، الجامعة}
  };
  % Uncomment to add logo (top-right corner):
  % \node at (\textwidth-3cm, -2.5cm) {\includegraphics[width=4cm]{logo.png}};
\end{tikzpicture}

\vspace{0.8cm}

% ============================================================
% POSTER BODY: three columns using minipage (landscape)
% ============================================================

\noindent
\begin{minipage}[t]{0.31\textwidth}

\posterblock{المقدمة}{%
% TODO: اكتب مقدمة البحث هنا.
هذا نص تجريبي للمقدمة. يصف المشكلة البحثية وأهميتها في مجال التقنية العصبية ورابط الدماغ والحاسوب.
}

\posterblock{الأهداف}{%
% TODO: اكتب أهداف البحث هنا.
\begin{enumerate}
    \item الهدف الأول من البحث
    \item الهدف الثاني من البحث
    \item الهدف الثالث من البحث
\end{enumerate}
}

\posterblock{المنهجية}{%
% TODO: اكتب المنهجية هنا.
\begin{itemize}
    \item الخطوة الأولى: \LR{Data Acquisition}
    \item الخطوة الثانية: المعالجة المسبقة
    \item الخطوة الثالثة: استخراج الميزات
    \item الخطوة الرابعة: التصنيف باستخدام \LR{Machine Learning}
\end{itemize}
}

\end{minipage}%
\hfill
\begin{minipage}[t]{0.31\textwidth}

\posterblock{النتائج}{%
% TODO: اكتب النتائج هنا.
\begin{center}
\begin{tabular}{|l|c|c|}
\hline
\textbf{الطريقة} & \textbf{الدقة} & \textbf{الزمن} \\
\hline
\LR{SVM} & 85.3\% & 12ms \\
\LR{CNN} & 92.1\% & 45ms \\
\LR{LSTM} & 89.7\% & 38ms \\
\hline
\end{tabular}
\end{center}
}

\posterblock{المناقشة}{%
% TODO: اكتب مناقشة النتائج هنا.
هذا نص تجريبي للمناقشة. يحلل النتائج ويقارنها بالأدبيات السابقة ويذكر أبرز النقاط.
}

\posterblock{الاستنتاجات}{%
% TODO: اكتب الاستنتاجات هنا.
هذا نص تجريبي للاستنتاجات. يلخص النتائج الرئيسية ويقترح أعمالاً مستقبلية.
}

\end{minipage}%
\hfill
\begin{minipage}[t]{0.31\textwidth}

\posterblock{الشكل التوضيحي}{%
% TODO: أضف شكل توضيحي هنا.
\begin{center}
% \includegraphics[width=\linewidth]{diagram.png}
\fbox{\parbox[c][5cm][c]{0.9\linewidth}{\centering مكان الشكل}}
\end{center}
}

\posterblock{المعادلات الرئيسية}{%
% TODO: أضف المعادلات هنا.
\begin{equation*}
y = f\left(\sum_{i=1}^{n} w_i x_i + b\right)
\end{equation*}
حيث $x_i$ المدخلات، $w_i$ الأوزان، $b$ الانحياز.
}

\posterblock{المراجع}{%
% TODO: اكتب المراجع هنا.
\begin{enumerate}
    \item Author, A. (2024). \emph{Title}. Journal, 1(1), 1--10.
    \item Author, B. (2023). \emph{Title}. Conference, pp. 1--5.
    \item Author, C. (2022). \emph{Title}. Book.
\end{enumerate}
}

\end{minipage}

\end{document}
